\section{Grundlagen und Stand der Technik}
\label{sec:grundlagen}

% \hl{structure notes start}
% \begin{itemize}
%     \item Aufzeigen dass es die Idee der Landing Zone bereits auf verschiedenen Plattformen gibt
%     \item Vergleich der verschiedenen Landing Zone Aspekte von den bekannten Cloud-Providern → Tabelle welche bereits im Vorhinein erstellt wurde
%     \item Kriterien aus bestehenden Arbeiten welche als Basis für Evaluationskriterien dienen können → Quasi aufzeigen dass die Kriterien welche wir definiert haben fundiert sind. Kriteren können zwar BIT/BV spezifische Sachen beinhalten sollten jedoch auch gewissermassen \flqq neutral\frqq\ sein
% \end{itemize}
% \hl{structure notes end}

Die Nutzung vom Begriff \ac{lz} selbst beschränkt sich in theoretischer Literatur mehrheitlich auf Referenzen von Cloud-Plattformen selbst und wird nicht explizit aufgegriffen. Ebenso finden sich keine ausführlichen Konzepte, wie \ac{lz} selbst zu strukturieren sind. In der Praxis bei den Cloud-Plattformen hingegen existiert ein genaueres Verständnis worum es sich bei einer \ac{lz} handelt. So wird sie nach \cite{microsoft_what_2025} bspw. bezeichnet als Möglichkeit zum Einrichten / Verwalten der Kundenumgebung mit dem Ziel eine Konsistenz bezüglich Sicherheit, Compliance, Betriebseffizienz etc. zu erhalten. Dieses Ziel kann als einer der treibenden Faktoren betrachtet werden, warum \ac{lz} relevant sind. Reduziert kann gesagt werden, dass durch eine \ac{lz} die Cloud-Plattformen als solche verfügbar und vor allem auf eine sichere Art nutzbar für eine Organisation gemacht wird \parencite[Kapitel 6]{mulder_multi-cloud_2020}. Eine sichere Nutzung bedingt zuerst dass die Veranwortlichkeiten zwischen Cloud-Plattform und dem Nutzer klar sind. Mit dem Modell der geteilten Verantwortung (oder auch Shared Responsibility Model) zeigen Plattformen deshalb auf, welchen Teil der Sicherung vom Kunden oder durch die Cloud-Plattform selbst erbracht werden müssen \parencite{aws_cloud_security_modell_nodate}. Für eine Organisationen heisst dies häufig, dass man sich bewusst werden muss welche Anforderungen man selbst hat. Diese können bspw. regulatorisch aus Dokumenten wie \ac{si001} wie auch organisatorisch durch Expertenteams / Plattformteams begründet sein. Die Themen welche relvant sind ähneln sich dabei über die Cloud-Plattformen hinweg. \ac{iam}, Network, Security sind dabei nur einige Beispiele.

Trotz der Gemeinsamkeiten dass eine Governance als notwendig verstanden wird, weisen die Cloud-Plattformen jedoch auch Unterschiede auf. Diese erstrecken sich von eigenen Metamodellen bis zu unterschiedlichen Implementation von Komponenten. Eine Organisation welche eine sichere Nutzung im multi-cloud Kontext sicherstellen möchte steht somit vor der Herausforderung, wie es die verschiedenen Cloud-Plattformen verwaltet. Hierfür existieren auf dem Markt bereits diverse Cloud Management Plattformen (bspw Flexera One, Cloud Bolt, HPE Morpheus), welche sich haupsächliche auf die operative Verwaltung und Bereitstellung von Cloud-Ressourcen spezialisieren. In anderen Worten wird addressiert wie die konkrete Implementation in einem multi-cloud Setup möglich ist. Offener dagegen ist wie die Strukturierung der Cloud-Ressourcen aussieht und welche Policies warum einer Cloud-Ressource zugewiesen werden müssen. Die Komplexität dieser Governance versuchen Cloud-Plattformen zwar individuell zu lösen, womit aber ein allgemeiner Ansatz fehlt. Diese Arbeit fokussiert sich deshalb auf einen Abstraktionsansatz für Governacne und nicht einer weiteren Cloud Management Plattform.

Für einen Abstraktionsansatz sind hierbei konkte die Security Frameworks relevant, dess Compliance sichergestellt werden soll. Ein prominentes Beispiel für die \ac{bv} ist \ac{si001}, welcher den IKT Grundschutz der \ac{bv} beschreibt. Ein anderes bekanntes Framework ist die \ac{nist} Special Publication 800-53. Beide genannten Regulation weisen mit Controls eine im Kern gleiche Struktur auf. Wichtiger jedoch ist, dass die korrekte Anwendung der Regulationen für Organisationen welche sie nutzen indirekt ein Ziel darstellt. Die Motvation kann dabei unterschiedlich sein (gesetztliche Vorgaben, selbstauferlegt etc.).

% start chat gpt
Mit den Regulation als Beispel stellt sich die Frage auf welcher Ebene eine solche Governance-Strukturierung erfolgen kann. Da Anforderungen aus regulatorischen Vorgaben oder organisationalen Zielsetzungen letztlich durch technische Architekturentscheidungen umgesetzt werden müssen, kann die Gestaltung von Landing Zones als ein architektonisches Problem verstanden werden.

Insbesondere im Multi-Cloud-Kontext entsteht dadurch die Notwendigkeit, eine Abstraktionsebene zu finden, welche unabhängig von spezifischen Implementationen einzelner Cloud-Provider funktioniert. Eine mögliche Perspektive hierfür bietet die Betrachtung von Fähigkeiten (\textit{Capabilities}) einer Organisation. Capabilities beschreiben stabilere, technologieunabhängige Möglichkeiten zur Erbringung bestimmter Leistungen, wie etwa das Management von Identitäten, die Segmentierung von Netzwerken oder die Überwachung von Systemzuständen.

Im Kontext von Enterprise-Architecture-Frameworks dienen Capabilities häufig als strukturierendes Element, um strategische Anforderungen mit konkreten Architekturartefakten in Beziehung zu setzen. Dadurch entsteht eine Grundlage, um Governance-Mechanismen wie Policies oder Controls systematisch einzuordnen und deren Umsetzung nachvollziehbar zu gestalten.
% end chat gpt
